%!TEX root = mainQlin.tex



% \begin{remark}
% \cnote{Discuss the runtime and space in case of $|\cA|$ is finite, or maximization over $\cA$ can be solved explicitly}.
% \end{remark}
% \begin{restatable}{lemma}{LEMunifcon}
% Suppose for any $\tau \in [k]$, $\x_\tau$ is sampled independently from distribution $\cD_\tau$, and $\bphi_\tau \in \cF_{\tau-1}$, where $\cF_\tau$ is the filtration until time $\tau$. Let $\Lambda = \sum_{\tau = 1}^k \bphi_\tau\bphi_\tau\trans + \lambda \I$, where $\lambda > 0$ is a regularization parameter. Then,  with high probability??, following holds for any $V \in \mathcal{V}$ simultaneously,
% \begin{equation*} 
% \norm{\sum_{\tau = 1}^k \bphi_\tau [V(x_\tau) - \E_{x \sim \cD_\tau} V(x)] }_{\Lambda^{-1}}
% \le \poly(d)
% \end{equation*}
% \end{restatable}

\section{Proof of Theorem \ref{thm:main_mis}} \label{app:proof_main_mis}
In this section, we prove Theorem \ref{thm:main_mis}. At a high level, the proof structure is similar to the structure in Appendix \ref{app:proof_main}. We will particularly focus on the parts that require different treatments in the misspecified setting.


\paragraph{Notation:} Throughout this section, we denote $\Lambda_h^k$, $\w^k_h$, and $Q_h^k$ as the parameters and the Q-value functions estimated in episode $k$. 
Denote the value function $V_h^k$ as $V_h^k(x) = \max_a Q_h^k(x, a)$. We denote $\pi_k$ as the greedy policy induced by $\{Q_h^k\}_{h=1}^H$.
To simplify the presentation, we denote $\bphi^k_h \defeq \bphi(x^{k}_h, a^{k}_h)$.

First, we establish a lemma that is the counterpart of Lemma \ref{prop:realizable} in the misspecified setting: for any policy $\pi$, its action-value function is always close to a linear function.

\begin{lemma}\label{lem:misspecified_error}
For a $\zeta$-nearly linear MDP, for any policy $\pi$, there exist corresponding weights $\{\w^{\pi}_h\}_{h\in[H]}$ where $\w^{\pi}_h = \btheta_h  + \int V^\pi_{h+1}(x') \dd \bmu_h (x')$ such that for any $(x, a, h) \in \cS\times\cA\times[H]$:
\begin{equation*}
|Q_h^\pi(x, a) - \la\bphi (x, a),  \w^{\pi}_h\ra| \le 2H\zeta.
\end{equation*}
% mispecified case, value, how close it is to the linear form.. 
% \cnote{fill this part.}
\end{lemma}
\begin{proof} Since $\bmu_h$ and $\btheta_h$ satisfy Eq.\eqref{eq:nearly_linear_transition}, we have:
\begin{align*}
&|Q^\pi_h(x, a) -  \la \bphi(x, a), \w^\pi_h\ra| \\
& \qquad \le  |r_h(x, a) - \la \bphi(x, a), \btheta_h\ra| + |\P_h \sind{V}{\pi}{h+1}(x, a) - \la \bphi(x, a),\int V^\pi_{h+1}(x') \dd \bmu_h (x') \ra| \\
& \qquad \le  \zeta + H\zeta \le 2H\zeta , 
\end{align*}
which finishes the proof.
\end{proof}

We can again show that the linear weights defined in Lemma \ref{lem:misspecified_error} are bounded.
% again show in the following two lemmas that the linear weights $\w_h$ that approximate the action-value functions or are in Algorithm \ref{algo:LSVI-UCB} are all bounded. 
\begin{lemma}[Bound on Weights of Value Functions] \label{lem:wn_policy_mis}
Under Assumption \ref{assumption:nearly_linear}, for any policy $\pi$, let $\{\w^{\pi}_h\}_{h\in[H]}$ be the corresponding weights as defined in Lemma \ref{lem:misspecified_error}.
% such that $\w^\pi_h = \btheta_h  + \int V^\pi_{h+1}(x') \dd \bmu_h (x')$ where $\bmu_h$ and $\btheta_h$ satisfy Eq.\eqref{eq:nearly_linear_transition}.
Then, we have
\begin{equation*}
\forall h \in [H], \quad \norm{\w^\pi_h} \le 2H\sqrt{d}.
\end{equation*}
\end{lemma}

\begin{proof}
Under the normalization conditions of Assumption \ref{assumption:nearly_linear}, the reward at each step is in $[0, 1]$, thus $V^\pi_{h+1}(x') \le H$ for any state $x'$. Therefore, $\norm{\btheta_h} \le \sqrt{d}$, and $\norm{\int V^\pi_{h+1}(x') \dd \bmu_h (x')} \le H\sqrt{d}$, which finishes the proof.
\end{proof}


Similar to Lemma \ref{lem:stochastic_term}, we also bound the stochastic noise in concentration.
\begin{lemma} \label{lem:stochastic_term_mis}
Under the setting of Theorem \ref{thm:main_mis}, let $c_{\beta}$ be the constant in our choice of $\beta_k$ (i.e. $\beta_k = c_\beta \cdot (d\sqrt{\logt}+\zeta \sqrt{kd})H$),
There exists an absolute constant $C$ that is independent of $c_{\beta}$ such that for any fixed $p\in[0, 1]$, if we let $\fE$ be the event that:
\begin{equation*}
\forall (k, h)\in [K]\times [H]: \quad \norm{\sum_{\tau = 1}^{k-1} \bphi^\tau_h [V^{k}_{h+1}(x^\tau_{h+1}) - \P_h V^{k}_{h+1}(x_h^\tau, a_h^\tau)]}_{(\Lambda^k_h)^{-1}}
\le C\cdot dH\sqrt{\chi}, 
\end{equation*}
where $\chi = \log [2(c_\beta+1)dT/p]$, then $\P(\fE) \ge 1-p/2$.
\end{lemma}
\begin{proof}
The proof is essentially the same as the proof for Lemma \ref{lem:stochastic_term}, with the only difference that $\beta_k$ is now bounded by $c_\beta (d\sqrt{\logt} +\zeta\sqrt{Kd})H$ instead of $c_\beta dH\sqrt{\logt}$ as in Lemma \ref{lem:stochastic_term}. Because $\zeta \le 1$ as in Assumption \ref{assumption:nearly_linear}, the new bound of $\beta_k$ only affects the choice of absolute $C$ in Lemma \ref{lem:stochastic_term_mis}. 
\end{proof}




% \begin{lemma}[Bound on Weights in Algorithm]\label{lem:wn_estimate_mis}
% In algorithm \ref{algo:LSVI-UCB}, we have for all $(k, h) \in [K] \times [H]$ that:
% \begin{equation*}
% \norm{\w^k_h} \le 2H \sqrt{dk/\lambda}.
% % 2Hk/\lambda.
% \end{equation*}
% \end{lemma}

% \begin{proof}
% For any vector $\v \in \R^d$, we have:
% \begin{align*}
% |\v\trans \w^k_h| & = |\v\trans (\Lambda^k_h)^{-1} \sum_{\tau=1}^{k-1} \bphi^\tau_h [r(x^{\tau}_h, a^{\tau}_h) + \max_a Q_{h+1}(x^{\tau}_{h+1}, a)]|\\
% & \le \sum_{\tau = 1}^{k-1}  |\v\trans (\Lambda^k_h)^{-1} \bphi^\tau_h| \cdot 2H 
% \le \sqrt{(\sum_{\tau = 1}^{k-1}  \v\trans (\Lambda^k_h)^{-1}\v) \cdot (\sum_{\tau = 1}^{k-1}  (\bphi^\tau_h)\trans (\Lambda^k_h)^{-1}\bphi^\tau_h)} \cdot 2H\\
% & \le 2H \norm{\v}\sqrt{dk/\lambda}
% \end{align*}
% where the last step is due to Lemma \ref{lem:basic_ineq}. The remaining of proof follows from the fact that 
% $\norm{\w^k_h} = \max_{\v:\norm{\v} = 1} |\v\trans \w^k_h|$.
% \end{proof}

% \begin{lemma} \label{lem:stochastic_term}
% Let $\fE$ be the event that for any $(k, h)\in [K]\times [H]$:
% \begin{equation*}
% \norm{\sum_{\tau = 1}^{k-1} \bphi^\tau_h [V^{k}_{h+1}(x^\tau_{h+1}) - \P_h V^{k}_{h+1}(x_h^\tau, a_h^\tau)]}_{(\Lambda^k_h)^{-1}}
% \le \tilde{\mathcal{O}}(dH)
% \end{equation*}
% Then, we have $\P(\fE) \ge 1-\delta$.
% \end{lemma}
% \begin{proof}
% This is immediate from Lemma \ref{lem:wn_estimate}, Lemma \ref{lem:self_norm_covering} and Lemma \ref{lem:covering_number}.
% \cnote{TODO: fill in more details.}
% \end{proof}


In the misspecified case, we also need to bound an error term where the noise can be potentially adversarial instead of stochastic. The adversarial noise is precisely due to model misspecification.
\begin{lemma}\label{lem:infinite_error}
Let $\{\epsilon_\tau\}$ be any sequence so that $|\epsilon_\tau| \le B$ for any $\tau$. Then, we have for any $(h, k) \in [H]\times[K]$ and any $\bphi \in \R^d$ that:
\begin{equation*}
\bigg |\bphi\trans (\Lambda^k_h)^{-1} \sum_{\tau = 1}^{k-1} \bphi^\tau_h \epsilon_\tau \bigg|
\le B\sqrt{dk\bphi\trans (\Lambda^k_h)^{-1}  \bphi}. 
\end{equation*}
\end{lemma}

\begin{proof}
By the Cauchy-Schwarz inequality,
\begin{align*}
\bigg |\bphi\trans (\Lambda^k_h)^{-1} \sum_{\tau = 1}^{k-1} \bphi^\tau_h \epsilon_\tau \bigg| \le& \sum_{\tau = 1}^{k-1}  |\bphi\trans (\Lambda^k_h)^{-1} \bphi^\tau_h| \cdot B
\le \sqrt{\biggl [ \sum_{\tau = 1}^{k-1}  \bphi\trans (\Lambda^k_h)^{-1}\bphi\biggr ]  \cdot \biggl [ \sum_{\tau = 1}^{k-1}  (\bphi^\tau_h)\trans (\Lambda^k_h)^{-1}\bphi^\tau_h\biggr ] } \cdot B\\
% & \le \norm{\Lambda_h^{-1}} \sum_{\tau=1}^k \norm{\bphi(x^{\tau}_h, a^{\tau}_h)} |r(x^{\tau}_h, a^{\tau}_h) + \max_a Q_{h+1}(x^{\tau}_{h+1}, a)|
\le& B \sqrt{dk\bphi\trans (\Lambda^k_h)^{-1}  \bphi},
\end{align*}
where the last inequality is due to Lemma \ref{lem:basic_ineq}.
\end{proof}

Now we are ready to prove the key lemma, which is the counterpart of Lemma \ref{lem:basic_relation}.

\begin{lemma}\label{lem:basic_relation_mis}There exists an absolute constant $c_\beta$ such that for $\beta_k = c_\beta \cdot (d\sqrt{\logt}+\zeta \sqrt{kd})H$ where $\logt = \log (2dT/p)$, and
for any fixed policy $\pi$, on the event $\fE$ defined in Lemma \ref{lem:stochastic_term_mis}, we have for all $(x, a, h, k) \in \cS\times\cA\times[H]\times[K]$ that:
\begin{equation*}
\la\bphi(x, a), \w^k_h\ra - Q_h^\pi(x, a)  =  \P_h (V^{k}_{h+1} - V^{\pi}_{h+1})(x, a) + \Delta^k_h(x, a),
\end{equation*}
for some $\Delta^k_h(x, a)$ that satisfies $|\Delta^k_h(x, a)| \le \beta_k \sqrt{\bphi(x, a)\trans (\Lambda^k_h)^{-1}  \bphi(x, a)} + 4H\zeta$.
\end{lemma}

\begin{proof}
By Lemma \ref{lem:misspecified_error}, there exists $\w^\pi_h = \btheta_h  + \int V^\pi_{h+1}(x') \dd \bmu_h (x')$ so that
for any $(x, a, h) \in \cS\times \cA \times [H]$:
\begin{equation*}
|Q_h^\pi(x, a) - \la\bphi (x, a),  \w^{\pi}_h\ra| \le 2H\zeta.
\end{equation*}
On the other hand, let $\tilde{\P}(\cdot|x, a) = \la\bphi(x, a), \bmu_h(\cdot)\ra$. Then, for any $(x, a) \in \cS \times \cA$, we have
\begin{equation*}
\la \bphi(x, a), \w^\pi_h\ra = \la \bphi(x, a), \btheta_h\ra + \tilde{\P}_h V^\pi_{h+1}(x, a).
\end{equation*}
This further gives
\begin{align*}
\lefteqn{\w^k_h -  \w^\pi_h
= (\Lambda^k_h)^{-1} \sum_{\tau = 1}^{k-1} \bphi^\tau_h [r^\tau_h + V^{k}_{h+1}(x^\tau_{h+1})]- \w^\pi_h}  \\
&= (\Lambda^k_h)^{-1} \bigg \{-\lambda \w^\pi_h + \sum_{\tau = 1}^{k-1} \bphi^\tau_h \bigl [r^{\tau}_h + V^{k}_{h+1}(x^\tau_{h+1}) - (\bphi^\tau_h)\trans\btheta_h - \tilde{\P}_h \sind{V}{\pi}{h+1}(x_h^\tau, a_h^\tau) \bigl ] \bigg\} \\
&= \underbrace{-\lambda (\Lambda^k_h)^{-1} \w^\pi_h}_{\q_1} + 
\underbrace{(\Lambda^k_h)^{-1} \sum_{\tau = 1}^{k-1} \bphi^\tau_h \bigl  [V^{k}_{h+1}(x^\tau_{h+1}) - \P_h V^{k}_{h+1}(x_h^\tau, a_h^\tau) \bigr ]}_{\q_2} 
 + \underbrace{(\Lambda^k_h)^{-1}\sum_{\tau = 1}^{k-1} \bphi^\tau_h \tilde{\P}_h (V^{k}_{h+1} - V^\pi_{h+1})(x_h^\tau, a_h^\tau)}_{\q_3}\\
  & \ \ + \underbrace{(\Lambda^k_h)^{-1}\sum_{\tau = 1}^{k-1} \bphi^\tau_h  \bigl [r^{\tau}_h - (\bphi^\tau_h)\trans\btheta_h + (\P_h - \tilde{\P}_h) V^k_{h+1}(x_h^\tau, a_h^\tau) \bigr ]}_{\q_4}. 
\end{align*}
Now, we bound the terms on the right-hand side individually.
For the first term,
\begin{equation*}
|\la \bphi(x, a), \q_1\ra| = |\lambda \la \bphi(x, a), (\Lambda^k_h)^{-1} \w^\pi_h\ra|
\le \sqrt{\lambda} \norm{\w_h^\pi} \sqrt{\bphi(x, a)\trans (\Lambda^k_h)^{-1}  \bphi(x, a)}. 
\end{equation*}
For the second term, given the event $\fE$ defined in Lemma \ref{lem:stochastic_term_mis}, we have:
\begin{equation*}
|\la \bphi(x, a), \q_2\ra| \le c_0 \cdot dH \sqrt{\chi} \sqrt{\bphi(x, a)\trans (\Lambda^k_h)^{-1}  \bphi(x, a)}, 
\end{equation*}
for an absolute constant $c_0$ independent of $c_\beta$, and $\chi = \log [2(c_\beta+1)dT/p]$. For the third term,
\begin{align*}
\la \bphi(x, a), \q_3\ra &= \la \bphi(x, a), (\Lambda^k_h)^{-1}\sum_{\tau = 1}^{k-1} \bphi^\tau_h \tilde{\P}_h (V^{k}_{h+1} - V^\pi_{h+1})(x_h^\tau, a_h^\tau)\ra\\
&= \la \bphi(x, a), (\Lambda^k_h)^{-1}\sum_{\tau = 1}^{k-1} \bphi^\tau_h (\bphi^\tau_h)\trans \int (V^{k}_{h+1} - V^\pi_{h+1})(x') \dd \bmu_h(x')\ra\\
&= \underbrace{\la \bphi(x, a), \int (V^{k}_{h+1} - V^\pi_{h+1})(x') \dd \bmu_h(x')\ra}_{p_1}
\underbrace{-\lambda \la \bphi(x, a), (\Lambda^k_h)^{-1}\int (V^{k}_{h+1} - V^\pi_{h+1})(x') \dd \bmu_h(x') \ra}_{p_2}, 
% &= \P_h (V^{k}_{h+1} - V^\star_{h+1})(x, a) -\lambda \la \bphi(x, a), (\Lambda^k_h)^{-1}\int (V^{k}_{h+1} - V^\star_{h+1})(x') \dd \bmu(x') \ra
\end{align*}
where by definition of $\tilde{\P}_h$, we have
\begin{align*}
p_1   = \tilde{\P}_h (V^{k}_{h+1} - V^\pi_{h+1})(x, a) , \qquad 
|p_2|  
% \ge \sqrt{\lambda} \norm{\int (V^{k}_{h+1} - V^\star_{h+1})(x') \dd \bmu(x')} \sqrt{\bphi(x, a)\trans (\Lambda^k_h)^{-1}  \bphi(x, a)}
\le 2 H \sqrt{d\lambda} \sqrt{\bphi(x, a)\trans (\Lambda^k_h)^{-1}  \bphi(x, a)}.
\end{align*}
Since $\|\tilde{\P}_h - \P_h\|_{\mathrm{TV}} \le \zeta$, we have 
\begin{equation*}
|p_1 - \P_h (V^{k}_{h+1} - V^\pi_{h+1})(x, a)| = |(\P_h-\tilde{P}_h) (V^{k}_{h+1} - V^\pi_{h+1})(x, a)| \le 2H\zeta.
\end{equation*}
For the fourth term, by Lemma \ref{lem:infinite_error}, we have
\begin{equation*}
|\la \bphi(x, a), \q_4\ra| \le 2H\zeta\sqrt{dk \bphi(x, a)\trans (\Lambda^k_h)^{-1}  \bphi(x, a)}. 
\end{equation*}

Finally, since $\la \bphi(x, a), \w^k_h - \w^\pi_h\ra =\la \bphi(x, a), \q_1 + \q_2 + \q_3 + \q_4\ra$, we have:
\begin{equation*}
|\la\bphi(x, a), \w^k_h\ra - Q_h^\pi(x, a)  -  \P_h (V^{k}_{h+1} - V^{\pi}_{h+1})(x, a)| \le (c' d\sqrt{\chi} + 2\zeta\sqrt{kd})H \sqrt{\bphi(x, a)\trans (\Lambda^k_h)^{-1}  \bphi(x, a)} + 4H\zeta,
\end{equation*}
for an absolute constant $c'$ independent of $c_{\beta}$. As in the proof of Lemma \ref{lem:basic_relation}, to prove this lemma, we only need to show that there exists a choice of absolute constant $c_\beta$ so that $c_\beta \ge 2$, and
\begin{equation*}
c' \sqrt{\iota + \log(c_\beta + 1)} \le c_\beta \sqrt{\iota}  \text{~for any~}\iota \in [\log 2, \infty).
\end{equation*}
This can be done by an picking absolute constant $c_\beta$ that satisfies $c'\sqrt{\log 2 + \log(c_\beta + 1)} \le c_\beta \sqrt{\log 2}$.
\end{proof}

Given Lemma \ref{lem:basic_relation_mis}, we can now easily proceed to prove that $Q^k_h$ is a upper bound of $Q^\star_h$ up to an error that depends linearly on the misspecification $\zeta$.

\begin{lemma}[UCB] \label{lem:UCB_misspecified}
Under the setting of Theorem \ref{thm:main_mis}, on the event $\fE$ defined in Lemma \ref{lem:stochastic_term_mis}, we have $Q^k_h(x, a) \ge Q^\star_h(x, a) - 4H(H+1-h)\zeta$ for all $(x, a, h, k) \in \cS\times\cA\times[H]\times[K]$.
\end{lemma}

\begin{proof}
% \cnote{Notes in plan.pdf, Be careful about initialization condition (needs to work out)}
We prove this lemma by induction. 

First, we consider the base case. The statement holds for the last step $H$, i.e., $Q^k_H(x, a) \ge Q^\star_H(x, a) - 4H\zeta$. Since the value function at $H+1$ step is zero, by Lemma \ref{lem:basic_relation_mis}, we have:
\begin{equation*}
|\la\bphi(x, a), \w^k_H\ra - Q_H^\star(x, a)| \le \beta_k \sqrt{\bphi(x, a)\trans (\Lambda^k_H)^{-1}  \bphi(x, a)} + 4H\zeta. 
\end{equation*}
Therefore, we obtain that 
\begin{equation*}
Q_H^\star(x, a) - 4H\zeta \le \min\{\la\bphi(x, a), \w^k_H\ra + \beta_k\sqrt{\bphi(x, a)\trans (\Lambda^k_H)^{-1}  \bphi(x, a)}, H\} = Q^k_H(x, a). 
\end{equation*}
Now, suppose the statement holds true at step $h+1$, and consider step $h$. Again, by Lemma \ref{lem:basic_relation_mis}, we have:
\begin{equation*}
|\la\bphi(x, a), \w^k_h\ra - Q_h^\star(x, a) -  \P_h (V^{k}_{h+1} - V^{\star}_{h+1})(x, a)| \le \beta_k \sqrt{\bphi(x, a)\trans (\Lambda^k_h)^{-1}  \bphi(x, a)}
+ 4H\zeta. 
\end{equation*}
By the induction assumption that $\P_h (V^{k}_{h+1} - V^{\star}_{h+1})(x, a) \ge -4H(H-h)\zeta$, we have:
\begin{equation*}
Q_h^\star(x, a) - 4H(H+1-h)\zeta \le \min\{\la\bphi(x, a), \w^k_h\ra + \beta_k\sqrt{\bphi(x, a)\trans (\Lambda^k_h)^{-1}  \bphi(x, a)}, H\} = Q^k_h(x, a).
\end{equation*}
Therefore, we conclude the proof of this lemma.
\end{proof}

The gap $\delta^k_h = V^k_h(x^{k}_{h}) - V^{\pi_k}_h(x^{k}_{h})$ also has a recursive formula similar to Lemma \ref{lem:recursive}.

\begin{lemma}[Recursive formula]\label{lem:recursive_mis}
Let $\delta^k_h = V^k_h(x^{k}_{h}) - V^{\pi_k}_h(x^{k}_{h})$, and $\zeta^k_{h+1} = 
\E[\delta^k_{h+1}|x^{k}_h, a^{k}_h] - \delta_{h+1}^k$. Then, on the event $\fE$ defined in Lemma \ref{lem:stochastic_term_mis}, we have the following for any $(k, h)\in [K]\times [H]$:
\begin{equation*}
% V^k_h(x^{k+1}_{h}) - V^{\pi_k}_h(x^{k+1}_{h})
% \le V^k_h(x^{k+1}_{h}) - V^{\pi_k}_h(x^{k+1}_{h})
\delta^k_h \le \delta^k_{h+1} + \zeta^k_{h+1} + 2\beta_k\sqrt{(\bphi^{k}_h)\trans (\Lambda^k_h)^{-1}  \bphi^{k}_h}.
\end{equation*}
% \cnote{Note this is correct because our analysis uses uniform convergence and is loose in $d$, so $\beta$ is large, verify the correctness of this lemma.}
\end{lemma}

\begin{proof} This is because by Lemma \ref{lem:basic_relation_mis}, we have for any $(x, a, h, k) \in \cS \times \cA \times [H] \times [K]$:
\begin{equation*}
Q^k_h(x, a) - Q^{\pi_k}_h(x, a)
\le \P_h (V^{k}_{h+1} - V^{\pi_k}_{h+1})(x, a)
+ 2\beta_k \sqrt{\bphi(x, a)\trans (\Lambda^k_h)^{-1}  \bphi(x, a)} + 4H\zeta. 
\end{equation*}
Finally, by Algorithm \ref{algo:LSVI-UCB} and the definition of $V^{\pi_k}$ we have
$$\delta^k_h = Q^k_h(x^{k}_h, a^{k}_h) - Q^{\pi_k}_h(x^{k}_h, a^{k}_h), $$
which finishes the proof.
\end{proof}

Finally, we are ready to combine all previous lemmas to prove the main theorem in the misspecified setting.

\THMmainmis*

\begin{proof}[Proof of Theorem \ref{thm:main_mis}]
	The proof of this theorem is similar to that of  Theorem \ref{thm:main}. 
	We condition on the event  $\fE$ defined in Lemma \ref{lem:stochastic_term_mis}.
	For for any $(k,h) \in [K]\times [H]$, 
	we define $\delta^k_h = V^k_h(x^{k}_{h}) - V^{\pi_k}_h(x^{k}_{h})$.   
	By Lemma \ref{lem:UCB_misspecified}, we have 
	$Q_1^k (x,a) \geq Q_1^*(x,a ) -4H^2  \zeta $ for all $k\in [K]$, which implies that $\sind{V}{\star}{1} (\ind{x}{k}{1}) - \sind{V}{\pi_k}{1} (\ind{x}{k}{1}) \leq \delta_1^k + 4H^2  \zeta$.
	Furthermore, by  Lemma \ref{lem:recursive_mis}, on the event $\fE$ we have:
\begin{align}\label{eq:mis1}
\text{Regret}(K) & =  \sum_{k=1}^{K} \left[\sind{V}{\star}{1} (\ind{x}{k}{1}) - \sind{V}{\pi_k}{1} (\ind{x}{k}{1})\right]
\le \sum_{k=1}^{K} [ \delta^k_1 + 4H^2 \zeta] \nn \\
& \le \sum_{k=1}^{K}\sum_{h=1}^H  \zeta^k_{h} + 2\sum_{k=1}^{K}\beta_k\sum_{h=1}^H \sqrt{(\bphi^{k}_h)\trans (\Lambda^k_h)^{-1}  \bphi^{k}_h}
+ 4HT\zeta,
\end{align}
where we use the fact that $T = HK$.
Since $\{ \zeta^k_{h} \}$  is a   martingale difference sequence with each term bounded by $2H$,
 the Azuma-Hoeffding inequality implies that 
 \begin{equation}\label{eq:mis2}
 \sum_{k=1}^{K}\sum_{h=1}^H  \zeta^k_{h}  \leq  \sqrt{2 T H^2 \cdot \log ( 2 / p)} \leq 2 H  \sqrt{ T\logt}
 \end{equation} 
 holds with  probability at least $1 -p / 2$, where $\logt = \log (2dT/p)$. 
 Moreover, by the Cauchy-Schwarz inequality, we have 
 \begin{equation}
 \label{eq:mis3}
  \sum_{k=1}^{K} \beta_k  \sqrt{(\bphi^{k}_h)\trans (\Lambda^k_h)^{-1}  \bphi^{k}_h}   \leq \biggl [ \sum_{k=1}^K \beta_k^2 \biggr ] ^{1/2} \cdot \biggl [ \sum_{k=1}^K (\bphi^{k}_h)\trans (\Lambda^k_h)^{-1}  \bphi^{k}_h \biggl] ^{1/2 }  . 
 \end{equation}
Similarly to Eq.~\eqref{eq:final08}, we have 
\begin{equation}
\label{eq:mis4}
 \biggl [ \sum_{k=1}^K (\bphi^{k}_h)\trans (\Lambda^k_h)^{-1}  \bphi^{k}_h \biggl] ^{1/2 }  \leq \sqrt{ 2 d \logt}. 
\end{equation}
 Moreover, since we set  $ \beta_k = c \cdot (d\sqrt{\logt} + \zeta\sqrt{kd})H  $ for some absolute constant $c>0$, we have 
 \begin{equation*}
 \sum_{k=1}^K \beta_k^2  =  c^2 \sum_{k=1}^K (d\sqrt{\logt} +  \zeta\sqrt{kd}) ^2  H^2  \leq 2 c^2 \sum_{k=1}^K  ( d^2\logt + \zeta^2 k d) H^2  \le 2c^2(d^2HT\logt + \zeta^2T^2d),
 \end{equation*}
 % & =    2 C^2 d^2 H^2 K \logt^2 + 2 C \zeta^2 d \cdot H^2  \cdot \iota ^2 \cdot  \bigg( \sum_{k=1}^K k \biggr ) \leq  2 C^2 d^2 H T \logt^2 + C \zeta^2 d \cdot T^2  \cdot \iota ^2 ,
 % \$
 which implies that 
 \begin{equation}
 \label{eq:mis5}
 \bigg( \sum_{k=1}^K \beta_k^2 \biggr )^{1/2} \leq \sqrt{2} c\bigl ( d \sqrt{H T \logt} + \zeta  T \sqrt{d} \bigr).
 \end{equation}
 Therefore, combining Eq.~\eqref{eq:mis3}, Eq.~\eqref{eq:mis4}, and Eq.~\eqref{eq:mis5}, we have 
 \#\label{eq:mis6}
 \sum_{k=1}^{K}\beta_k\sum_{h=1}^H \sqrt{(\bphi^{k}_h)\trans (\Lambda^k_h)^{-1}  \bphi^{k}_h}
  % \le c' H \cdot \sqrt{2d \logt } \cdot \bigl (d \sqrt{H T \logt} + \zeta T \sqrt{d} \bigr )
   \le 2c\cdot (\sqrt{d^3 H^3 T\logt^2}  + \zeta d  H T\sqrt{\logt} ).
 \#
  Finally, combining Eq.~\eqref{eq:mis1}, Eq.~\eqref{eq:mis2}, and  Eq.~\eqref{eq:mis6}, we obtain 
  \$
  \text{Regret}(K)  \le c''\cdot (\sqrt{d^3 H^3 T\logt^2}  + \zeta d  H T\sqrt{\logt} ),
  \$
 for some absolute constant $c''$. This concludes the proof of the theorem.
 \end{proof}
%  \iffalse
%  $ \beta_k = \Theta( (d + \zeta\sqrt{kd})H \cdot \iota),$ 
% where the last step is because (1) the first term is martingale; (2) the second term is standard in linear bandit (matrix/vector pigeon-hole).

% $$ \le \tilde{\mathcal{O}}(\sqrt{d^3 H^3 T} + \zeta dHT)$$
% \cnote{TODO: fill in more details.}
% \fi 
